\documentclass[11pt,oneside]{book}

% ---------------------------------------------------------------
% Packages
% ---------------------------------------------------------------
\usepackage[utf8]{inputenc}
\usepackage[T1]{fontenc}
\usepackage[english]{babel}
\usepackage{amsmath,amssymb,amsthm,mathtools}
\usepackage{geometry}
\usepackage{titlesec}
\usepackage{fancyhdr}
\usepackage{xcolor}
\usepackage{enumitem}
\usepackage{titling}
\usepackage[bookmarks,hidelinks,pdfusetitle]{hyperref}

\geometry{paperwidth=7in, paperheight=10in, top=20mm, bottom=22mm, inner=20mm, outer=16mm}

\definecolor{inkblue}{RGB}{25,35,60}
\definecolor{rule}{RGB}{120,120,120}

% Chapter / section styling
\titleformat{\chapter}[display]
  {\normalfont\bfseries\color{inkblue}}
  {\Large\MakeUppercase{\chaptertitlename} \thechapter}
  {6pt}
  {\Huge}
\titlespacing*{\chapter}{0pt}{10pt}{28pt}

\titleformat{\section}
  {\normalfont\Large\bfseries\color{inkblue}}{\thesection}{0.8em}{}
\titleformat{\subsection}
  {\normalfont\large\bfseries\color{inkblue}}{\thesubsection}{0.8em}{}

\pagestyle{fancy}
\fancyhf{}
\renewcommand{\headrulewidth}{0.3pt}
\fancyhead[LE,RO]{\thepage}
\fancyhead[RE]{\nouppercase{\leftmark}}
\fancyhead[LO]{\nouppercase{\rightmark}}

\setlength{\parindent}{1.2em}
\setlength{\parskip}{0.2em}
\tolerance=1200
\emergencystretch=3em
\hbadness=10000

\title{Unified Configuration}
\author{Kirill Nikitenko}
\date{2026}

\begin{document}

% ---------------------------------------------------------------
% Cover / Title page
% ---------------------------------------------------------------
\begin{titlepage}
\centering
\vspace*{2.2cm}

{\color{rule}\rule{0.55\textwidth}{0.6pt}}

\vspace{1.6cm}
{\Large\scshape\color{inkblue} A Research Framework in Theoretical Physics}\\[0.8cm]

{\fontsize{30}{34}\selectfont\bfseries\color{inkblue} Unified\\[0.15cm] Configuration}\\[0.6cm]

{\Large\itshape Work Field Theory}

\vspace{1.6cm}
{\color{rule}\rule{0.55\textwidth}{0.6pt}}

\vfill

{\large Kirill Nikitenko}\\[0.3cm]
{\large 2026}

\vspace{1.2cm}
\end{titlepage}

\frontmatter

% ---------------------------------------------------------------
% Abstract
% ---------------------------------------------------------------
\chapter*{Abstract}
\addcontentsline{toc}{chapter}{Abstract}

Modern physics describes nature through several successful but conceptually
disjoint languages: the states and operators of quantum theory, and the
curved spacetime of general relativity. \emph{Unified Configuration}
develops an alternative, exploratory layer of description in which the
primitive object of physics is neither a particle nor a field in isolation,
but a \emph{configuration} --- a stable, distributed state of energy in an
abstract configuration space $\mathcal{C}$.

Within this framework, matter is characterized as a local minimum of an
energy functional $\mathcal{E}[q]$; work is reinterpreted as motion of a
configuration through this landscape; and interactions, radiation, and
gravitation are treated as different manifestations of the geometry and
curvature of the same underlying space. Stability, transition, and decay are
analyzed through the Hessian of the energy functional, and physical law is
recast as a statement about the accessibility and connectivity of
configuration space rather than about isolated equations of motion.

The book develops this language systematically, from foundational
definitions through a mathematical formalism, into applications spanning
particle stability, field structure, and cosmological evolution, and closes
by proposing configuration space itself as an object of long-term
scientific and engineering inquiry. The framework is offered as a
structure for generating testable models rather than as a finished theory;
its scope and epistemic status are set out explicitly in the chapter that
follows.

% ---------------------------------------------------------------
% How to Read This Work
% ---------------------------------------------------------------
\chapter*{How to Read This Work}
\addcontentsline{toc}{chapter}{How to Read This Work}

This chapter states, in advance, what kind of document follows and what
kind of claims it makes. Three points define its scope.

\bigskip
\noindent\textbf{1. This is a hypothesis, not a replacement.}
Quantum mechanics, quantum field theory, and general relativity remain the
empirically established descriptions of nature within their tested domains,
and nothing in this work is intended to contradict or supersede them.
\emph{Unified Configuration} is a research hypothesis proposing a different
\emph{level} of description --- a configuration-space language from which
existing theories are conjectured to emerge as limiting cases. Every
definition, equation, and interpretation offered here should be read as a
proposal under active construction, not as an established result.

\bigskip
\noindent\textbf{2. Experiment governs the model, not the reverse.}
Where a prediction of this framework conflicts with observation, the
appropriate response is to revise the specific part of the model
responsible for that prediction --- never to dismiss, reinterpret, or set
aside the data. Declarative language in the chapters that follow (``is,''
``requires,'' ``therefore'') is the working form in which a hypothesis is
stated clearly; it is not a claim of settled fact. When evidence narrows or
contradicts a particular construction, that construction is to be treated
as provisional and open to revision, without this implying that the
surrounding architecture is thereby invalidated.

\bigskip
\noindent\textbf{3. The goal is a structure, not a final answer.}
This book does not aim to deliver a closed set of results. Its purpose is
to provide a common structure --- configurations, energy functionals, and
stability conditions --- within which specific, falsifiable models can be
built, tested, and revised. Progress, in this sense, is measured less by
whether any individual equation survives contact with experiment, and more
by whether the framework continues to generate questions that can be
checked against the world.

\bigskip
\noindent Readers are encouraged to keep these three points in mind
throughout, particularly in later chapters where the language becomes more
concrete and the temptation to read a working hypothesis as a finished
theory is correspondingly greater.

\tableofcontents

\mainmatter

% ---------------------------------------------------------------
% Body
% ---------------------------------------------------------------
\input{body_reformatted.tex}

\end{document}
